% --- Archon physics formalization source begin ---
% archon:physics
% archon:covers PhyXMiniProblems/problem_phyx_mini_0076.lean
% archon:source-report reports/phyx_mini/problem_phyx_mini_0076.source.json

\chapter{Physics problem phyx\_mini\_0076}
\label{ch:PhyXMiniProblems_problem_phyx_mini_0076}

\paragraph{Problem source.}
\#\# Question

In the figure, at what distance are are parallel rays from the left focused to a point?

\#\# Answer choices

- A: A: 14cm

- B: B: 26cm

- C: C: 15cm

- D: D: 19.5cm

\#\# Figure caption (auxiliary; use the image as primary evidence)

The image depicts a diagram of a lens system. It includes:

- **Objects**:
  - Two lenses: a diverging lens on the left with focal length \textbackslash{}( f\_1 = -10 \textbackslash{}, \textbackslash{}text\{cm\} \textbackslash{}) and a converging lens on the right with focal length \textbackslash{}( f\_2 = 10 \textbackslash{}, \textbackslash{}text\{cm\} \textbackslash{}).
  - A series of parallel purple arrows entering from the left, indicating light rays.

- **Relationships**:
  - The lenses are positioned 20 cm apart.
  - The light enters the diverging lens and then converges towards the converging lens.

- **Text** notations indicating the focal lengths of the lenses and the distance between them:
  - \textbackslash{}( f\_1 = -10 \textbackslash{}, \textbackslash{}text\{cm\} \textbackslash{})
  - \textbackslash{}( f\_2 = 10 \textbackslash{}, \textbackslash{}text\{cm\} \textbackslash{})
  - A horizontal line marks the optical axis.

This setup illustrates a scenario in optics involving a pair of lenses with specified focal lengths and positions.

\#\# Dataset metadata

Geometrical Optics; Physical Model Grounding Reasoning, Multi-Formula Reasoning

\paragraph{Recorded answer/context.}
C: C: 15cm

\paragraph{Figure/image path.}
/root/proposal\_for\_physic/science-mango/phyx\_mini\_run/phyx\_data/test\_image/76.png

\paragraph{Formalization target.}
create a compiling Lean file with sorry bodies at `PhyXMiniProblems/problem\_phyx\_mini\_0076.lean`. The Lean declarations must preserve the physical quantities, dimensions or dimensional roles, figure labels, governing-law hypotheses, and final relation expressed by this problem.
Use Mathlib/Physlib names found through LeanExplore where available. If a physics API is missing, introduce faithful local abstractions rather than scalar placeholder aliases.

\begin{theorem}[Physics formalization target]
\label{thm:physics:phyx_mini_0076:target}
\lean{PhyXMiniProblems.ProblemPhyXMini0076.problem_phyx_mini_0076}
\uses{def:physics:phyx-mini-0076:phyxminiproblems-problemphyxmini0076-lensmirrorsetup, def:physics:phyx-mini-0076:phyxminiproblems-problemphyxmini0076-mirrorimagedistanceincentimeters, def:physics:phyx-mini-0076:phyxminiproblems-problemphyxmini0076-hasphysicallensmirrorconfiguration, def:physics:phyx-mini-0076:phyxminiproblems-problemphyxmini0076-matchesfigurereadouts, def:physics:phyx-mini-0076:phyxminiproblems-problemphyxmini0076-satisfiesparaxiallensmirrorlaws, def:physics:phyx-mini-0076:phyxminiproblems-problemphyxmini0076-matchesanswer}
The assigned autoformalize agent should translate this physics problem into Lean declarations in the covered file, with theorem and lemma proof bodies written as `by sorry`.
\end{theorem}
\begin{proof}
This is an autoformalization task, not a proof task. Produce statements that can later be proved by the physics prover without weakening the physical model.
\end{proof}

% --- Archon declaration topology begin ---
\section{Declaration topology}
\begin{definition}[Signed Length Quantity]
  \label{def:physics:phyx-mini-0076:phyxminiproblems-problemphyxmini0076-signedlengthquantity}
  \lean{PhyXMiniProblems.ProblemPhyXMini0076.SignedLengthQuantity}
  A signed physical length, independent of the unit chosen to read it
\end{definition}
\begin{proof}
  This is the declared model definition of Signed Length Quantity.
\end{proof}
\begin{definition}[signed Length Readout]
  \label{def:physics:phyx-mini-0076:phyxminiproblems-problemphyxmini0076-signedlengthreadout}
  \lean{PhyXMiniProblems.ProblemPhyXMini0076.signedLengthReadout}
  \uses{def:physics:phyx-mini-0076:phyxminiproblems-problemphyxmini0076-signedlengthquantity}
  The scalar readout of a signed physical length in a specified unit system
\end{definition}
\begin{proof}
  This is the declared model definition of signed Length Readout.
\end{proof}
\begin{definition}[centimeter Unit Choices]
  \label{def:physics:phyx-mini-0076:phyxminiproblems-problemphyxmini0076-centimeterunitchoices}
  \lean{PhyXMiniProblems.ProblemPhyXMini0076.centimeterUnitChoices}
  Unit choices with centimeters selected as the length unit
\end{definition}
\begin{proof}
  This is the declared model definition of centimeter Unit Choices.
\end{proof}
\begin{definition}[signed Length In Centimeters]
  \label{def:physics:phyx-mini-0076:phyxminiproblems-problemphyxmini0076-signedlengthincentimeters}
  \lean{PhyXMiniProblems.ProblemPhyXMini0076.signedLengthInCentimeters}
  \uses{def:physics:phyx-mini-0076:phyxminiproblems-problemphyxmini0076-signedlengthquantity, def:physics:phyx-mini-0076:phyxminiproblems-problemphyxmini0076-signedlengthreadout, def:physics:phyx-mini-0076:phyxminiproblems-problemphyxmini0076-centimeterunitchoices}
  The signed centimeter readout of an axial length or position
\end{definition}
\begin{proof}
  This is the declared model definition of signed Length In Centimeters.
\end{proof}
\begin{definition}[Optical Axis Point]
  \label{def:physics:phyx-mini-0076:phyxminiproblems-problemphyxmini0076-opticalaxispoint}
  \lean{PhyXMiniProblems.ProblemPhyXMini0076.OpticalAxisPoint}
  Distinguished points on the horizontal principal axis in the figure
\end{definition}
\begin{proof}
  This is the declared model definition of Optical Axis Point.
\end{proof}
\begin{definition}[Thin Lens Kind]
  \label{def:physics:phyx-mini-0076:phyxminiproblems-problemphyxmini0076-thinlenskind}
  \lean{PhyXMiniProblems.ProblemPhyXMini0076.ThinLensKind}
  The two qualitative kinds of thin lens
\end{definition}
\begin{proof}
  This is the declared model definition of Thin Lens Kind.
\end{proof}
\begin{definition}[Spherical Mirror Kind]
  \label{def:physics:phyx-mini-0076:phyxminiproblems-problemphyxmini0076-sphericalmirrorkind}
  \lean{PhyXMiniProblems.ProblemPhyXMini0076.SphericalMirrorKind}
  The two qualitative kinds of spherical mirror
\end{definition}
\begin{proof}
  This is the declared model definition of Spherical Mirror Kind.
\end{proof}
\begin{definition}[Axial Propagation Direction]
  \label{def:physics:phyx-mini-0076:phyxminiproblems-problemphyxmini0076-axialpropagationdirection}
  \lean{PhyXMiniProblems.ProblemPhyXMini0076.AxialPropagationDirection}
  Propagation direction along the oriented horizontal optical axis
\end{definition}
\begin{proof}
  This is the declared model definition of Axial Propagation Direction.
\end{proof}
\begin{definition}[Ray Bundle Profile]
  \label{def:physics:phyx-mini-0076:phyxminiproblems-problemphyxmini0076-raybundleprofile}
  \lean{PhyXMiniProblems.ProblemPhyXMini0076.RayBundleProfile}
  Angular profile of a paraxial ray bundle relative to the principal axis
\end{definition}
\begin{proof}
  This is the declared model definition of Ray Bundle Profile.
\end{proof}
\begin{definition}[Incident Ray Bundle]
  \label{def:physics:phyx-mini-0076:phyxminiproblems-problemphyxmini0076-incidentraybundle}
  \lean{PhyXMiniProblems.ProblemPhyXMini0076.IncidentRayBundle}
  \uses{def:physics:phyx-mini-0076:phyxminiproblems-problemphyxmini0076-axialpropagationdirection, def:physics:phyx-mini-0076:phyxminiproblems-problemphyxmini0076-raybundleprofile}
  Qualitative information carried by the incident collection of rays
\end{definition}
\begin{proof}
  This is the declared model definition of Incident Ray Bundle.
\end{proof}
\begin{definition}[Optical Approximation]
  \label{def:physics:phyx-mini-0076:phyxminiproblems-problemphyxmini0076-opticalapproximation}
  \lean{PhyXMiniProblems.ProblemPhyXMini0076.OpticalApproximation}
  The optical approximation in which the imaging laws are asserted
\end{definition}
\begin{proof}
  This is the declared model definition of Optical Approximation.
\end{proof}
\begin{definition}[Thin Lens]
  \label{def:physics:phyx-mini-0076:phyxminiproblems-problemphyxmini0076-thinlens}
  \lean{PhyXMiniProblems.ProblemPhyXMini0076.ThinLens}
  \uses{def:physics:phyx-mini-0076:phyxminiproblems-problemphyxmini0076-signedlengthquantity, def:physics:phyx-mini-0076:phyxminiproblems-problemphyxmini0076-opticalaxispoint, def:physics:phyx-mini-0076:phyxminiproblems-problemphyxmini0076-thinlenskind}
  A thin lens with a signed focal length and a center on the optical axis
\end{definition}
\begin{proof}
  This is the declared model definition of Thin Lens.
\end{proof}
\begin{definition}[Spherical Mirror]
  \label{def:physics:phyx-mini-0076:phyxminiproblems-problemphyxmini0076-sphericalmirror}
  \lean{PhyXMiniProblems.ProblemPhyXMini0076.SphericalMirror}
  \uses{def:physics:phyx-mini-0076:phyxminiproblems-problemphyxmini0076-signedlengthquantity, def:physics:phyx-mini-0076:phyxminiproblems-problemphyxmini0076-opticalaxispoint, def:physics:phyx-mini-0076:phyxminiproblems-problemphyxmini0076-sphericalmirrorkind}
  A spherical mirror with a signed focal length and an axial vertex
\end{definition}
\begin{proof}
  This is the declared model definition of Spherical Mirror.
\end{proof}
\begin{definition}[Lens Mirror Setup]
  \label{def:physics:phyx-mini-0076:phyxminiproblems-problemphyxmini0076-lensmirrorsetup}
  \lean{PhyXMiniProblems.ProblemPhyXMini0076.LensMirrorSetup}
  \uses{def:physics:phyx-mini-0076:phyxminiproblems-problemphyxmini0076-signedlengthquantity, def:physics:phyx-mini-0076:phyxminiproblems-problemphyxmini0076-opticalaxispoint, def:physics:phyx-mini-0076:phyxminiproblems-problemphyxmini0076-incidentraybundle, def:physics:phyx-mini-0076:phyxminiproblems-problemphyxmini0076-opticalapproximation, def:physics:phyx-mini-0076:phyxminiproblems-problemphyxmini0076-thinlens, def:physics:phyx-mini-0076:phyxminiproblems-problemphyxmini0076-sphericalmirror}
  The physical components, ray input, and labeled axial points of the source figure. The positions of the virtual source and requested final focus remain unknown dimensionful quantities in this structure
\end{definition}
\begin{proof}
  This is the declared model definition of Lens Mirror Setup.
\end{proof}
\begin{definition}[axial Position Readout]
  \label{def:physics:phyx-mini-0076:phyxminiproblems-problemphyxmini0076-axialpositionreadout}
  \lean{PhyXMiniProblems.ProblemPhyXMini0076.axialPositionReadout}
  \uses{def:physics:phyx-mini-0076:phyxminiproblems-problemphyxmini0076-signedlengthreadout, def:physics:phyx-mini-0076:phyxminiproblems-problemphyxmini0076-opticalaxispoint, def:physics:phyx-mini-0076:phyxminiproblems-problemphyxmini0076-lensmirrorsetup}
  The coordinate readout of an axial label in an arbitrary unit system
\end{definition}
\begin{proof}
  This is the declared model definition of axial Position Readout.
\end{proof}
\begin{definition}[lens Focal Length Readout]
  \label{def:physics:phyx-mini-0076:phyxminiproblems-problemphyxmini0076-lensfocallengthreadout}
  \lean{PhyXMiniProblems.ProblemPhyXMini0076.lensFocalLengthReadout}
  \uses{def:physics:phyx-mini-0076:phyxminiproblems-problemphyxmini0076-signedlengthreadout, def:physics:phyx-mini-0076:phyxminiproblems-problemphyxmini0076-lensmirrorsetup}
  The signed focal-length readout of the diverging lens
\end{definition}
\begin{proof}
  This is the declared model definition of lens Focal Length Readout.
\end{proof}
\begin{definition}[mirror Focal Length Readout]
  \label{def:physics:phyx-mini-0076:phyxminiproblems-problemphyxmini0076-mirrorfocallengthreadout}
  \lean{PhyXMiniProblems.ProblemPhyXMini0076.mirrorFocalLengthReadout}
  \uses{def:physics:phyx-mini-0076:phyxminiproblems-problemphyxmini0076-signedlengthreadout, def:physics:phyx-mini-0076:phyxminiproblems-problemphyxmini0076-lensmirrorsetup}
  The signed focal-length readout of the spherical mirror
\end{definition}
\begin{proof}
  This is the declared model definition of mirror Focal Length Readout.
\end{proof}
\begin{definition}[mirror Object Distance Readout]
  \label{def:physics:phyx-mini-0076:phyxminiproblems-problemphyxmini0076-mirrorobjectdistancereadout}
  \lean{PhyXMiniProblems.ProblemPhyXMini0076.mirrorObjectDistanceReadout}
  \uses{def:physics:phyx-mini-0076:phyxminiproblems-problemphyxmini0076-lensmirrorsetup, def:physics:phyx-mini-0076:phyxminiproblems-problemphyxmini0076-axialpositionreadout}
  The mirror's object distance: mirror vertex minus the virtual-source coordinate. It is positive for the real object seen by the pictured mirror
\end{definition}
\begin{proof}
  This is the declared model definition of mirror Object Distance Readout.
\end{proof}
\begin{definition}[mirror Image Distance Readout]
  \label{def:physics:phyx-mini-0076:phyxminiproblems-problemphyxmini0076-mirrorimagedistancereadout}
  \lean{PhyXMiniProblems.ProblemPhyXMini0076.mirrorImageDistanceReadout}
  \uses{def:physics:phyx-mini-0076:phyxminiproblems-problemphyxmini0076-lensmirrorsetup, def:physics:phyx-mini-0076:phyxminiproblems-problemphyxmini0076-axialpositionreadout}
  The mirror-to-focus distance in front of the mirror. Since the axis points to the right, this is mirror vertex minus final-focus coordinate
\end{definition}
\begin{proof}
  This is the declared model definition of mirror Image Distance Readout.
\end{proof}
\begin{definition}[mirror Object Distance In Centimeters]
  \label{def:physics:phyx-mini-0076:phyxminiproblems-problemphyxmini0076-mirrorobjectdistanceincentimeters}
  \lean{PhyXMiniProblems.ProblemPhyXMini0076.mirrorObjectDistanceInCentimeters}
  \uses{def:physics:phyx-mini-0076:phyxminiproblems-problemphyxmini0076-centimeterunitchoices, def:physics:phyx-mini-0076:phyxminiproblems-problemphyxmini0076-lensmirrorsetup, def:physics:phyx-mini-0076:phyxminiproblems-problemphyxmini0076-mirrorobjectdistancereadout}
  The mirror's physical object-distance readout in centimeters
\end{definition}
\begin{proof}
  This is the declared model definition of mirror Object Distance In Centimeters.
\end{proof}
\begin{definition}[mirror Image Distance In Centimeters]
  \label{def:physics:phyx-mini-0076:phyxminiproblems-problemphyxmini0076-mirrorimagedistanceincentimeters}
  \lean{PhyXMiniProblems.ProblemPhyXMini0076.mirrorImageDistanceInCentimeters}
  \uses{def:physics:phyx-mini-0076:phyxminiproblems-problemphyxmini0076-centimeterunitchoices, def:physics:phyx-mini-0076:phyxminiproblems-problemphyxmini0076-lensmirrorsetup, def:physics:phyx-mini-0076:phyxminiproblems-problemphyxmini0076-mirrorimagedistancereadout}
  The requested mirror-to-focus distance, read in centimeters
\end{definition}
\begin{proof}
  This is the declared model definition of mirror Image Distance In Centimeters.
\end{proof}
\begin{definition}[Has Focal Signs Consistent With Kinds]
  \label{def:physics:phyx-mini-0076:phyxminiproblems-problemphyxmini0076-hasfocalsignsconsistentwithkinds}
  \lean{PhyXMiniProblems.ProblemPhyXMini0076.HasFocalSignsConsistentWithKinds}
  \uses{def:physics:phyx-mini-0076:phyxminiproblems-problemphyxmini0076-centimeterunitchoices, def:physics:phyx-mini-0076:phyxminiproblems-problemphyxmini0076-lensmirrorsetup, def:physics:phyx-mini-0076:phyxminiproblems-problemphyxmini0076-lensfocallengthreadout, def:physics:phyx-mini-0076:phyxminiproblems-problemphyxmini0076-mirrorfocallengthreadout}
  The Gaussian focal-length sign agrees with the optical component kind
\end{definition}
\begin{proof}
  This is the declared model definition of Has Focal Signs Consistent With Kinds.
\end{proof}
\begin{definition}[Has Physical Lens Mirror Configuration]
  \label{def:physics:phyx-mini-0076:phyxminiproblems-problemphyxmini0076-hasphysicallensmirrorconfiguration}
  \lean{PhyXMiniProblems.ProblemPhyXMini0076.HasPhysicalLensMirrorConfiguration}
  \uses{def:physics:phyx-mini-0076:phyxminiproblems-problemphyxmini0076-centimeterunitchoices, def:physics:phyx-mini-0076:phyxminiproblems-problemphyxmini0076-lensmirrorsetup, def:physics:phyx-mini-0076:phyxminiproblems-problemphyxmini0076-axialpositionreadout, def:physics:phyx-mini-0076:phyxminiproblems-problemphyxmini0076-hasfocalsignsconsistentwithkinds}
  Qualitative component and ray roles supplied by the primary image. No position or distance for the requested final focus is included here
\end{definition}
\begin{proof}
  This is the declared model definition of Has Physical Lens Mirror Configuration.
\end{proof}
\begin{definition}[Matches Figure Readouts]
  \label{def:physics:phyx-mini-0076:phyxminiproblems-problemphyxmini0076-matchesfigurereadouts}
  \lean{PhyXMiniProblems.ProblemPhyXMini0076.MatchesFigureReadouts}
  \uses{def:physics:phyx-mini-0076:phyxminiproblems-problemphyxmini0076-centimeterunitchoices, def:physics:phyx-mini-0076:phyxminiproblems-problemphyxmini0076-lensmirrorsetup, def:physics:phyx-mini-0076:phyxminiproblems-problemphyxmini0076-axialpositionreadout, def:physics:phyx-mini-0076:phyxminiproblems-problemphyxmini0076-lensfocallengthreadout, def:physics:phyx-mini-0076:phyxminiproblems-problemphyxmini0076-mirrorfocallengthreadout}
  Numerical and graphical readouts in the primary image: f₁ = -10 cm, f₂ = 10 cm, a 20 cm component separation, and the principal axis. The unknown virtual-source and final-focus positions are deliberately absent
\end{definition}
\begin{proof}
  This is the declared model definition of Matches Figure Readouts.
\end{proof}
\begin{definition}[Satisfies Paraxial Lens Mirror Laws]
  \label{def:physics:phyx-mini-0076:phyxminiproblems-problemphyxmini0076-satisfiesparaxiallensmirrorlaws}
  \lean{PhyXMiniProblems.ProblemPhyXMini0076.SatisfiesParaxialLensMirrorLaws}
  \uses{def:physics:phyx-mini-0076:phyxminiproblems-problemphyxmini0076-lensmirrorsetup, def:physics:phyx-mini-0076:phyxminiproblems-problemphyxmini0076-axialpositionreadout, def:physics:phyx-mini-0076:phyxminiproblems-problemphyxmini0076-lensfocallengthreadout, def:physics:phyx-mini-0076:phyxminiproblems-problemphyxmini0076-mirrorfocallengthreadout, def:physics:phyx-mini-0076:phyxminiproblems-problemphyxmini0076-mirrorobjectdistancereadout, def:physics:phyx-mini-0076:phyxminiproblems-problemphyxmini0076-mirrorimagedistancereadout}
  The governing paraxial laws for the lens--mirror sequence. Parallel rays leaving a diverging lens behave as if emitted from the lens's incident-side focal point. The mirror law is the Gaussian equation 1/f = 1/p + 1/q, written in the division-free, dimensionally homogeneous form f * (p + q) = p * q. Both laws are required in every unit system and neither supplies the requested 15 cm result
\end{definition}
\begin{proof}
  This is the declared model definition of Satisfies Paraxial Lens Mirror Laws.
\end{proof}
\begin{definition}[Answer Choice]
  \label{def:physics:phyx-mini-0076:phyxminiproblems-problemphyxmini0076-answerchoice}
  \lean{PhyXMiniProblems.ProblemPhyXMini0076.AnswerChoice}
  Labels of the four displayed multiple-choice answers
\end{definition}
\begin{proof}
  This is the declared model definition of Answer Choice.
\end{proof}
\begin{definition}[distance In Centimeters]
  \label{def:physics:phyx-mini-0076:phyxminiproblems-problemphyxmini0076-answerchoice-distanceincentimeters}
  \lean{PhyXMiniProblems.ProblemPhyXMini0076.AnswerChoice.distanceInCentimeters}
  \uses{def:physics:phyx-mini-0076:phyxminiproblems-problemphyxmini0076-answerchoice}
  The focus distance in centimeters printed beside each answer choice
\end{definition}
\begin{proof}
  This is the declared model definition of distance In Centimeters.
\end{proof}
\begin{definition}[Matches Answer]
  \label{def:physics:phyx-mini-0076:phyxminiproblems-problemphyxmini0076-matchesanswer}
  \lean{PhyXMiniProblems.ProblemPhyXMini0076.MatchesAnswer}
  \uses{def:physics:phyx-mini-0076:phyxminiproblems-problemphyxmini0076-lensmirrorsetup, def:physics:phyx-mini-0076:phyxminiproblems-problemphyxmini0076-mirrorimagedistanceincentimeters, def:physics:phyx-mini-0076:phyxminiproblems-problemphyxmini0076-answerchoice}
  Exact agreement between the derived focus distance and a displayed choice
\end{definition}
\begin{proof}
  This is the declared model definition of Matches Answer.
\end{proof}
\begin{lemma}[mirror Object Distance eq thirty]
  \label{lem:physics:phyx-mini-0076:phyxminiproblems-problemphyxmini0076-mirrorobjectdistance-eq-thirty}
  \lean{PhyXMiniProblems.ProblemPhyXMini0076.mirrorObjectDistance_eq_thirty}
  \uses{def:physics:phyx-mini-0076:phyxminiproblems-problemphyxmini0076-lensmirrorsetup, def:physics:phyx-mini-0076:phyxminiproblems-problemphyxmini0076-mirrorobjectdistanceincentimeters, def:physics:phyx-mini-0076:phyxminiproblems-problemphyxmini0076-matchesfigurereadouts, def:physics:phyx-mini-0076:phyxminiproblems-problemphyxmini0076-satisfiesparaxiallensmirrorlaws}
  The lens's virtual source is 10 cm left of the lens and hence 30 cm in front of the mirror
\end{lemma}
\begin{proof}
  The conclusion follows from the governing relations and figure readouts named in the statement of mirror Object Distance eq thirty.
\end{proof}
% --- Archon declaration topology end ---
% --- Archon physics formalization source end ---
